\section{Tabu Search: improving one route plan}
TS stores a route plan directly. It starts from CW. If CW has too many routes, the code
removes the smallest routes and reinserts their clinics into $K$ remaining slots. This can
temporarily overload a route. The search then repairs and improves the plan.

\subsection*{Moves and constant-time cost changes}
\textbf{Relocate} moves a clinic to another position or route. \textbf{Swap} exchanges two clinics.
\textbf{2-opt} reverses a segment within one route. \textbf{2-opt*} exchanges tails of two routes.
The code supports all four. Candidate lists focus the search on nearby clinics, with $k=20$
in the main experiments.

For a relocate between distinct routes, remove $u$ from edges $(a,u),(u,b)$ and insert it between
$c$ and $d$. The change in cost is
\[
\begin{split}
 \Delta={}&d(a,b)-d(a,u)-d(u,b)\\
          &+d(c,u)+d(u,d)-d(c,d).
\end{split}
\]
Only six edge costs are needed, regardless of the number of clinics. Cached route loads give
the overload change in constant time too. Applying the chosen move can take longer because
route arrays and positions must be updated.

\subsection*{Tabu memory, aspiration and overload}
Removed undirected edges become tabu for a random 10--25 iterations. A move reintroducing
one is normally forbidden. \emph{Aspiration} allows it if it creates a new best feasible solution.

Moves are ranked by cost plus $\lambda$ times total overload. If the new plan is overloaded,
multiply $\lambda$ by 1.3; otherwise divide it by 1.3, within fixed bounds. This discourages
long periods of overload while letting the search move between clinic assignments. Initially,
$\lambda$ is average depot distance divided by average demand. Only zero-overload plans with
at most $K$ routes can become saved answers. TS follows the memory-based search idea in [7].

\begin{algorithm}[h]\caption{Tabu Search (\texttt{tabu.cpp}, moves in \texttt{routes.hpp})}
\begin{algorithmic}[1]
\State Start from CW, reducing to at most $K$ route slots if necessary
\State Initialise tabu expiry times, penalty and best feasible answer
\While{time remains}
 \State Enumerate moves; compute their cost and overload changes
 \State Reject tabu moves unless aspiration holds
 \State Choose the best admissible penalised move, even if current cost worsens
 \State Apply it and mark removed edges tabu
 \State Adapt the overload penalty
 \State Save a feasible plan if it improves the best answer
 \State After stagnation, restart from the best answer with random relocations
\EndWhile
\State Return the best feasible answer
\end{algorithmic}\end{algorithm}

\takeaway{GA searches a population of orders. TS changes an explicit route plan, using short-term
memory to discourage cycling back to recently reversed choices.}
